\subsection{Full Leaderboard}

Table~\ref{tab:leaderboard} reports every model, every metric, sorted
by MAE. Bold rows are the two best models, which come out of
Section~\ref{sec:models}'s SoTA push (PatchTST seed ensemble alone, and
the mean ensemble of PatchTST and LSTM). Italic rows are the naive
baselines.

\begin{table}[!htbp]
\centering
\caption{Full leaderboard --- 20 models, 6 metrics. Mean~$\pm$~standard
deviation across the four rolling-origin windows. Lower is better for
every metric except R\textsuperscript{2}. Bold: the two best models
overall (Section~\ref{sec:models}, SoTA push). Italic: naive baselines.
Foundation models are zero-shot (no training on this dataset).}
\label{tab:leaderboard}
\setlength{\tabcolsep}{4pt}
\scriptsize
\resizebox{\textwidth}{!}{%
\begin{tabular}{@{}lrrrrrr@{}}
\toprule
Model & MAE (kWh) & RMSE (kWh) & MAPE & MASE & Bias & R\textsuperscript{2} \\
\midrule
\textbf{Ensemble PatchTST + LSTM}        & \textbf{25.09\,$\pm$\,9.4}  & \textbf{28.88\,$\pm$\,11.3}  & 58.5\% & \textbf{0.63} & $-$2.42  & $-$0.13 \\
\textbf{PatchTST (3-seed ensemble)}      & \textbf{25.19\,$\pm$\,8.9}  & 30.74\,$\pm$\,11.9            & 58.7\% & 0.64          & $-$1.92  & $-$0.32 \\
ARIMAX (2,1,3)                            & 25.52\,$\pm$\,14.4          & 30.51\,$\pm$\,15.6            & 53.9\% & 0.65          & $-$8.08  & $-$0.19 \\
LSTM seq14\_h64\_l2                       & 25.68\,$\pm$\,10.6          & 28.47\,$\pm$\,11.2            & 59.4\% & 0.65          & $-$2.93  & $-$0.12 \\
SARIMAX (2,1,4)$\times$(1,1,2,7)          & 25.89\,$\pm$\,12.8          & 31.12\,$\pm$\,13.7            & 54.0\% & 0.65          & $-$7.96  & $-$0.27 \\
TimesFM 2.0 -- Google (500M)              & 26.23\,$\pm$\,12.6          & 30.79\,$\pm$\,14.6            & 53.5\% & 0.66          & $-$9.71  & $-$0.22 \\
ARIMA (2,1,3)                             & 26.31\,$\pm$\,12.4          & 30.78\,$\pm$\,13.7            & 58.3\% & 0.67          & $-$5.82  & $-$0.25 \\
TimesFM 1.0 -- Google (200M)              & 26.32\,$\pm$\,11.1          & 30.60\,$\pm$\,12.8            & 57.5\% & 0.67          & $-$6.87  & $-$0.25 \\
AutoARIMA non-seas. (0,1,1)               & 26.40\,$\pm$\,11.5          & 30.64\,$\pm$\,13.0            & 57.9\% & 0.67          & $-$6.39  & $-$0.25 \\
ARIMA (3,1,3)                             & 26.68\,$\pm$\,11.2          & 31.28\,$\pm$\,13.3            & 58.2\% & 0.67          & $-$6.19  & $-$0.30 \\
ARMA (3,0,3)                              & 26.71\,$\pm$\,12.3          & 30.75\,$\pm$\,13.9            & 58.7\% & 0.68          & $-$6.63  & $-$0.24 \\
SARIMA (3,1,4)$\times$(0,1,2,7)           & 26.86\,$\pm$\,11.9          & 31.92\,$\pm$\,13.5            & 57.7\% & 0.68          & $-$6.13  & $-$0.36 \\
Chronos-T5 base -- Amazon (200M)          & 27.31\,$\pm$\,11.3          & 32.75\,$\pm$\,14.1            & 58.7\% & 0.69          & $-$8.93  & $-$0.44 \\
AutoARIMA seas. (0,1,1)$\times$(1,0,1,7)  & 27.44\,$\pm$\,11.7          & 31.44\,$\pm$\,13.5            & 57.3\% & 0.69          & $-$8.82  & $-$0.31 \\
Chronos-T5 large -- Amazon (710M)         & 27.94\,$\pm$\,13.0          & 33.24\,$\pm$\,16.0            & 53.6\% & 0.71          & $-$14.26 & $-$0.42 \\
\midrule
Stacking NNLS (LOWO)                      & 26.08\,$\pm$\,9.8           & 28.99\,$\pm$\,10.8            & 61.0\% & 0.66          & $-$2.32  & $-$0.17 \\
Stacking Ridge (LOWO)                     & 32.29\,$\pm$\,8.1           & 37.22\,$\pm$\,8.3             & 77.7\% & 0.82          & $-$2.73  & $-$1.22 \\
\midrule
\textit{Historical mean}                  & \textit{30.13\,$\pm$\,2.2}  & \textit{34.87\,$\pm$\,3.0}    & \textit{89.9\%} & \textit{0.76} & \textit{+14.42} & \textit{$-$1.21} \\
\textit{Last value}                       & \textit{31.22\,$\pm$\,12.7} & \textit{37.08\,$\pm$\,13.8}   & \textit{61.0\%} & \textit{0.79} & \textit{$-$10.33} & \textit{$-$1.02} \\
\textit{Drift}                            & \textit{31.24\,$\pm$\,12.7} & \textit{37.09\,$\pm$\,13.8}   & \textit{61.1\%} & \textit{0.79} & \textit{$-$10.29} & \textit{$-$1.03} \\
\textit{ARIMA naive-seasonal (3,1,3)}     & \textit{34.74\,$\pm$\,12.1} & \textit{41.36\,$\pm$\,11.0}   & \textit{76.0\%} & \textit{0.88} & \textit{$-$4.95}  & \textit{$-$1.69} \\
\textit{Seasonal naive}                   & \textit{34.98\,$\pm$\,13.1} & \textit{41.57\,$\pm$\,12.4}   & \textit{79.6\%} & \textit{0.88} & \textit{$-$4.10}  & \textit{$-$1.68} \\
\bottomrule
\end{tabular}}
\end{table}

\subsection{Leaderboard Plot}

Figure~\ref{fig:leaderboard} visualises the same ranking but with the
error bars (one standard deviation across windows). Three observations
stand out:
\begin{enumerate}[leftmargin=*]
    \item The top of the leaderboard is tightly packed: every model
    in the top six is within 1~kWh MAE of the best, which is smaller
    than the standard deviation of every individual model.
    \item The Ensemble PatchTST + LSTM does \emph{not} have the
    smallest standard deviation --- that honour goes to the
    historical-mean baseline, which is trivially stable. The smallest
    std among forecasting models belongs to PatchTST alone
    (8.90~vs.\ 14.38 for ARIMAX): the seed-ensemble buys us almost
    40\% variance reduction over the previous best.
    \item All forecasting models comfortably beat the naive
    baselines. The MAE gap between the best model (25.09~kWh) and
    the best baseline (30.13~kWh) is $\sim$17\%, which is the
    learnable signal in the series.
\end{enumerate}

\begin{figure}[!htbp]
\centering
\includegraphics[width=0.95\textwidth]{fig_leaderboard}
\caption{Final leaderboard. Bars are MAE in kWh (lower is better);
error bars are one standard deviation across the four rolling-origin
windows. The Ensemble PatchTST + LSTM (top) is the best model
overall; the PatchTST seed-ensemble has the lowest std among
forecasting models.}
\label{fig:leaderboard}
\end{figure}

\subsection{Per-Window Analysis}

Figure~\ref{fig:per-window} breaks down the MAE by window for the top
six models. The takeaway here is window~3 (Nov 14--20, 2023), which
includes Thanksgiving: every model has its worst week here, with MAE
shooting up to 36--43~kWh. PatchTST is the one that loses the least
on that week (36.4~kWh vs.\ 38.8~kWh for the LSTM and $\sim$41~kWh
for ARIMAX), and that is the main reason it wins the leaderboard.
Outside window~3, the top models are essentially tied.

\begin{figure}[!htbp]
\centering
\includegraphics[width=0.95\textwidth]{fig_per_window}
\caption{Per-window MAE for the top six models. Window~3 (Thanksgiving
week) is hard for every model; the ranking on the other three windows
is much tighter. PatchTST loses less on the hard week, which is what
puts it (and its ensemble with LSTM) on top of the leaderboard.}
\label{fig:per-window}
\end{figure}

\subsection{Google vs.\ Amazon: Bigger Is Not Better}

The foundation-model rows of Table~\ref{tab:leaderboard} contain a
small surprise. Both Google TimesFM models clearly beat both Amazon
Chronos models (TimesFM 2.0: 26.23~kWh; Chronos-T5 large: 27.94~kWh,
the worst foundation model in our leaderboard). Inside each family
the size ordering is also informative: TimesFM 2.0 (500M) is only
marginally better than TimesFM 1.0 (200M), and Chronos-T5 large
(710M) is actually \emph{worse} than its smaller 200M sibling. Two
take-aways:
\begin{itemize}[leftmargin=*]
    \item \textbf{Foundation models do not win automatically.} The
    best foundation model (TimesFM 2.0) is at position 6 of our
    leaderboard, still beaten by ARIMAX, LSTM and SARIMAX in MAE.
    Zero-shot is a useful baseline but not a substitute for training
    on the target data.
    \item \textbf{More parameters means more prior, not more
    flexibility.} If the pre-training corpus is dominated by series
    that do not look like Boulder EV demand (random-walk-like
    financial series, for instance), a larger model amplifies the
    wrong prior. TimesFM transfers better because its pre-training
    mix includes more energy and load series; Chronos is trained on
    a different corpus and the larger checkpoint inherits more of
    that mismatch.
\end{itemize}

\subsection{What We Beat, and What We Did Not}

The honest claim, summarised:
\begin{itemize}[leftmargin=*]
    \item \textbf{Inside our protocol}, the Ensemble PatchTST + LSTM
    beats the strongest non-transformer model (ARIMAX) on both the
    mean MAE (25.09~vs.\ 25.52~kWh) and the standard deviation
    (9.36~vs.\ 14.38~kWh, ~35\% smaller). That is a real,
    apples-to-apples improvement.
    \item \textbf{Against the published literature}, no claim can be
    made in either direction. Koohfar~\cite{koohfar2023transformer},
    Wu~\cite{wu2025catformer} and
    Kyriakopoulos~\cite{kyriakopoulos2025comparative} all report
    metrics in normalised units (z-score or MinMax), on different
    station selections, and mostly on different horizons. Our 25~kWh
    and their 0.7--0.9 (normalised) are not in the same unit system.
\end{itemize}

\subsection{Stacking Did Not Work --- and Why}

For completeness we also tried two stacked meta-learners (Ridge and
NNLS, both leave-one-window-out, both trained on the six base learners
listed in Section~\ref{sec:models}). The results are in the middle
rows of Table~\ref{tab:leaderboard}: NNLS lands at 26.08~kWh MAE
(worse than the simple mean), and Ridge collapses to 32.29~kWh, which
is below several baselines. The reason is structural: 21 training
points per fold is not enough to fit reliable meta-weights. The Ridge
solution oscillates wildly between folds (we observed both very
negative and very large positive coefficients depending on which
window is held out), and NNLS concentrates almost all its weight on
the LSTM, throwing away the diversity the meta-learner was supposed
to exploit.

The dumbest combination wins: the unweighted mean of the two best
deep models, picked a priori. The lesson, which we re-state in
Section~\ref{sec:conclusions}, is that stacking needs a real
validation set --- not just a tiny held-out test set --- to learn
weights that generalise.
